% RESEARCH_DRAFT_SHA256: 97a8519ba79c5b7163531b8cc46ba3f9d1a97debb456f0f17686549d3e76a498
% Layout and source-reference wrappers expanded for the web renderer.
\documentclass[11pt]{article}
\usepackage[margin=25mm]{geometry}
\usepackage{fontspec,unicode-math}
\setmainfont{Latin Modern Roman}
\setmathfont{Latin Modern Math}
\usepackage{amsmath,amssymb,mathtools,xurl,hyperref}
\hypersetup{colorlinks=true,urlcolor=blue,linkcolor=blue}
\setcounter{secnumdepth}{0}
\setlength{\parindent}{0pt}
\setlength{\parskip}{5pt}
\setlength{\emergencystretch}{3em}
\newcommand{\sref}[2]{\hyperlink{source:#1:#2}{[S#2]}}
\newcommand{\cataloguescope}[1]{\paragraph{Source record.}#1}
\newcommand{\floor}[1]{\left\lfloor#1\right\rfloor}
\begin{document}
\section*{Research-attempt provenance}
This standalone research attempt retains the complete original definition body below,
including its historical source and preparation notes. Those notes describe the
earlier catalogue preparation; the new attempt and its verification are identified
separately after that body. No original ID, catalogue status, or source statement is changed.

% UnsolvedMath ID 30001569; source catalogue code OWR-4426-006
\section{Polyhedral Information Derived from Known Symmetries}\label{30001569}
{\small UnsolvedMath ID 30001569 · Computer Science}
\subsection{Definitions and mathematical statement}

\paragraph{Shared notation.}Standard number systems are $\mathbb N,\mathbb Z,\mathbb Q,\mathbb R,\mathbb C$. Specialized notation and index ranges follow the cited source and the definitions below.


A symmetry of a polyhedron maps it onto itself and therefore permutes its vertices and faces. A group orbit consists of all images of one point or face under the group action.
\paragraph{Record scope and qualifications.}
The catalogue title is retained as metadata. Named certificates, contracts, bounds or frameworks in that title are not asserted as results of the original source extract.
The source requests potentially useful information, without a definite invariant or algorithmic output requirement.

Source-record limitation: The request for any information derived from symmetry does not specify a mathematical proposition.
\paragraph{Statement recorded in the supplied source.}
Given a bounded polyhedron

$$P=\operatorname{conv}\{x_1,\ldots,x_k\}\subset\mathbb R^n$$

and its known symmetry group, what structural or algorithmic information can those symmetries provide about $P$? \href{https://doi.org/10.4171/owr/2010/38}{[S2]}
\paragraph{Source status.}
The supplied catalogue labels this source record partially solved; that metadata was checked against the supplied snapshot on 2026-10-11. The label is preserved for provenance and is not an independent verification of present open status or mathematical correctness. The catalogue literature assessment is dated 2026-08-21.
\subsection{Short English statement}
Use known polyhedral symmetries to derive useful structure or algorithms.
\subsection{Sources}
\begin{itemize}
\item[S1] Ulam AI and the UnsolvedMath Contributors, UnsolvedMath: A Curated Collection of Open Mathematics Problems, record 30001569 (OWR-4426-006). Catalogue curation and metadata: CC BY 4.0. The source extract is retained; definitions, English summaries and qualifications are editorial additions.
\url{https://huggingface.co/datasets/ulamai/UnsolvedMath}.
\item[S2] Mini-Workshop: Exploiting Symmetry in Optimization (2011), Oberwolfach Reports, DOI 10.4171/owr/2010/38, pp. 2245--2281
\url{https://doi.org/10.4171/owr/2010/38}.
\item[S3] Bremner et al., Computing symmetry groups of polyhedra
\url{https://doi.org/10.1112/S1461157014000400}.
\item Source-specific Oberwolfach Reports licences, including CC BY-SA 4.0 where applicable, are retained; see THIRD\_PARTY\_NOTICES.md.
\end{itemize}
\paragraph{Preparation and review.}
The original source and UnsolvedMath contributors retain attribution for the source extract. Codex (GPT-6; exact serving revision not exposed) assisted with definitions, English summaries, scope notes and formatting. No new human mathematical authorship or independently verified solution is asserted. Permanent source ID 30001569 and code OWR-4426-006 checked against the complete supplied registry and current repository. Distinct upstream records, including repeated questions and broad or damaged extracts, are retained for complete coverage of the source CS category. A repeated record is not a new independent conjecture. Complete mathematical-equivalence and current-literature review has not been performed. Preparation and final review follow contribution prompts C and D, with source gaps disclosed. Structural and rendering checks do not establish a new theorem or open status.
\label{end:30001569}

\section{Literature and context check}
The primary report on printed pages 2270--2271 already states that orbit barycenters describe the fixed-point polytope and explains the difficulty of restricting integer solutions to it. The following rederives that theorem directly, proves its optimization implication, and gives exact examples showing what it does not determine. This is source-context analysis, not a new symmetry theorem.
\section{Research attempt: Fixed-point projection and information lost by averaging}

\paragraph{Averaging gives the exact fixed-point polytope}
Let a finite group $G$ of affine transformations preserve a bounded polytope $P=\operatorname{conv}\{v_1,\ldots,v_k\}$. Define the affine averaging map
\[
 R(x)=\frac1{|G|}\sum_{g\in G}g(x).
\]
For $x\in P$, every $gx$ lies in $P$, so convexity gives $R(x)\in P$. For any $h\in G$, affine maps preserve averages and left multiplication permutes $G$, hence $hR(x)=R(x)$. Conversely if $x$ is fixed by every $g$, then $R(x)=x$. Therefore
\[
 R(P)=P\cap\operatorname{Fix}(G).
\]
An affine map preserves convex combinations, giving $R(P)=\operatorname{conv}\{R(v_i)\}$. On a vertex orbit each orbit point occurs $|\operatorname{Stab}(v_i)|$ times in the group sum, so $R(v_i)$ is exactly that orbit's barycenter. One barycenter per orbit thus determines the fixed-point polytope.

For a linear objective $c^Tx$ invariant under the action on $P$, averaging also preserves its value. Compactness of $P$ gives an optimum, whose average is a fixed optimum. Hence optimizing this objective over $P\cap\operatorname{Fix}(G)$ has the same optimum as over $P$. This conclusion is for the real convex problem.

\paragraph{Why integer optimization cannot simply take fixed points}
Let $P=\operatorname{conv}\{(1,0),(0,1)\}$ and let $G$ swap coordinates. Its integer feasible points are the two endpoints, while its fixed-point polytope is the singleton $(1/2,1/2)$ and contains no integer point. The constant invariant objective has feasible integer optima in $P$ but no feasible point after restricting integer variables to the fixed subspace. Thus exact real averaging does not preserve integrality even when all transformations preserve the integer lattice.

\paragraph{The projection also loses geometric information}
The square $P_1=[-1,1]^2$ and the diamond $P_2=\operatorname{conv}\{(2,0),(0,2),(-2,0),(0,-2)\}$ are both preserved by the same group generated by a quarter-turn and reflection. Both have one orbit of vertices with barycenter zero and fixed-point set $\{0\}$. Their areas are four and eight, respectively: the square is side two, and the diamond is the union of four right triangles of legs two. Their lattice-point counts are nine and thirteen: the square has three choices per coordinate, while the diamond has one point at Manhattan norm zero, four at norm one and eight at norm two. Their edge directions differ as well; for example $(1,1)$ is on both boundaries, whereas $(2,0)$ belongs only to the diamond. Thus the group and its barycenter projection do not reconstruct the full polytope, its volume or its integral fibers.

These examples do not deny that the original vertex list plus group contains enough data to compute such properties by other methods. They show that the tempting compression to fixed points alone discards information needed by those methods.

\section{Verification and adversarial audit}
The averaging proof covers affine actions and explicitly uses finiteness and convexity. Orbit-stabilizer multiplicity is counted before replacing group averages by orbit averages. The integer example has a lattice-preserving symmetry but a nonintegral average. The geometric examples use the same concrete linear symmetry matrices and different areas.
\paragraph{Audit corrections.} Marked the barycenter theorem as already in the primary source and retained the distinction between convex optimization and integer feasibility.
\paragraph{External computation.} None; symbolic arguments and boundary checks only.
\section{Scope and remaining gaps}
The broad request for useful structural or algorithmic information is not a single theorem to be completed. The known fixed-point result and its exact limitations are proved, but no general efficient recovery of the lost fibers or integer optimization method is supplied.
\paragraph{Integration review.} Codex (GPT-6; exact serving revision not exposed) reviewed the written arguments and expanded layout and source-reference wrappers. Repository metadata and the production web renderer were checked using Python 3.13 and JavaScript with MathJax 3.2.2. This model-assisted review is not independent human proof verification; no PDF compilation or formal proof checking is asserted.

\section{Completion Estimate}
\noindent\textbf{COMPLETION: 0\%}
The calculations below the preserved source statement investigate a concrete route,
but establish no substantiated new progress toward the entire original target.
Known special cases and failed sufficient conditions are not credited as a solution.
\section{References}
The original source references and catalogue attribution are preserved in the
Sources section above. The following URLs were actually consulted for this attempt;
retention of an additional original reference does not claim fresh consultation of it.
\begin{itemize}
\item \url{https://ems.press/content/serial-article-files/46297} (consulted 2026-10-11). 
\end{itemize}
\section{Model and contributor disclosure}
The parent Codex GPT-6 model has reviewed the written mathematical arguments at their declared scopes. This is model-assisted review, not independent human verification.
OpenAI gpt-6-astra, requested reasoning configuration ultra, generated this research
attempt on 2026-10-11. The exact serving revision is not exposed. Codex GPT-6
performed orchestration and repository-format review; its exact serving revision
is also not exposed. The mathematical route, boundary checks and final adversarial
self-audit followed contribution prompts A then B in
\texttt{docs/contribution\_prompts.md}, after reading \texttt{CONTRIBUTING.md}.
Python 3.14.2 assembled the standalone documents; any mathematical computation is
specified in Verification. No theorem prover was used. No human mathematical
authorship, substantial human proof correction, or independent human proof checking
is asserted. Model-assisted review, source retrieval, and structural checks do not
establish correctness or novelty. Full private reasoning transcripts are unavailable
as reproducibility artifacts; the complete arguments, inputs and outcomes actually
used are stated here or in the companion batch audit.
\section{Final status}
Unresolved as a new research contribution. This is a concrete route investigation,
not a claim of new complete resolution or a current-status certification. A known
answer reproduced for context does not become a new contribution. The catalogue's
source-status label is preserved; this attempt status does not assert that a
source-labelled solved question is open.
\end{document}
